\section{Lines in Three-Dimensional Space}

\begin{corebox}[The same construction in one more coordinate]
A line in space is still determined by one point and one nonzero direction. The
third coordinate changes the arithmetic, but not the geometric idea.
\end{corebox}

\begin{definitionbox}
Let
\[
P_0=(x_0,y_0,z_0),
\qquad
\mathbf v=(a,b,c)\neq\mathbf0.
\]
The line through $P_0$ with direction $\mathbf v$ is
\[
\boxed{L=\{P_0+t\mathbf v:t\in\mathbb R\}}.
\]
Equivalently,
\[
\boxed{
\begin{cases}
x=x_0+at,\\
y=y_0+bt,\\
z=z_0+ct.
\end{cases}
}
\]
\end{definitionbox}

\begin{interpretationbox}[Meaning of the parameter]
As $t$ varies, the point moves along the line. Positive and negative values give
the two orientations. Rescaling the direction vector changes only the speed at
which the parameter traverses the same set of points.
\end{interpretationbox}

\begin{center}
\begin{tikzpicture}[scale=0.9]
    \draw[axis] (0,0) -- (5.2,0) node[right] {$x$};
    \draw[axis] (0,0) -- (-1.7,-1.5) node[below left] {$y$};
    \draw[axis] (0,0) -- (0,4.3) node[above] {$z$};

    \coordinate (P) at (1.2,1.0);
    \coordinate (Q) at (4.2,3.0);
    \draw[subset] (0.5,0.53) -- (4.8,3.4) node[above right] {$L$};
    \fill[geometricSubsetColor] (P) circle (2.2pt) node[below left] {$P_0$};
    \draw[->,thick,geometricSubsetColor] (P) -- (Q) node[midway,above left] {$\mathbf v$};
\end{tikzpicture}
\end{center}

\subsection{The Line Through Two Points in Space}

\begin{derivationbox}[The direction is the displacement between the points]
Let
\[
P=(x_1,y_1,z_1),
\qquad
Q=(x_2,y_2,z_2)
\]
be distinct. Then
\[
Q-P=(x_2-x_1,\,y_2-y_1,\,z_2-z_1)
\]
is a nonzero direction vector.
\end{derivationbox}

\begin{theorembox}
The unique line through distinct points $P,Q\in\mathbb R^3$ is
\[
\boxed{L=\{P+t(Q-P):t\in\mathbb R\}}.
\]
In coordinates,
\[
\boxed{
\begin{cases}
x=x_1+t(x_2-x_1),\\
y=y_1+t(y_2-y_1),\\
z=z_1+t(z_2-z_1).
\end{cases}
}
\]
\end{theorembox}

\begin{examplebox}
For
\[
P=(1,0,2),
\qquad
Q=(3,4,1),
\]
a direction vector is $Q-P=(2,4,-1)$. Therefore
\[
(x,y,z)=(1,0,2)+t(2,4,-1),
\]
or
\[
\begin{cases}
x=1+2t,\\
y=4t,\\
z=2-t.
\end{cases}
\]
\end{examplebox}

\subsection{Symmetric Form}

\begin{derivationbox}[Eliminating the parameter]
If $a,b,c$ are all nonzero, then
\[
t=\frac{x-x_0}{a}
=\frac{y-y_0}{b}
=\frac{z-z_0}{c}.
\]
Hence
\[
\boxed{
\frac{x-x_0}{a}
=
\frac{y-y_0}{b}
=
\frac{z-z_0}{c}
}.
\]
\end{derivationbox}

\begin{remarkbox}
If a direction component is zero, the corresponding coordinate is constant. For
example, $c=0$ gives $z=z_0$. The parametric form handles this without division
by zero and is therefore the more fundamental representation.
\end{remarkbox}

\subsection{Parallel, Perpendicular, Intersecting, and Skew Lines}

Consider
\[
L_1=P_1+t\mathbf v_1,
\qquad
L_2=P_2+s\mathbf v_2.
\]

\begin{theorembox}
The directions are parallel precisely when
\[
\boxed{\mathbf v_1\times\mathbf v_2=\mathbf0}.
\]
They are perpendicular precisely when
\[
\boxed{\mathbf v_1\cdot\mathbf v_2=0}.
\]
\end{theorembox}

\begin{remarkbox}
Parallel directions do not by themselves distinguish identical lines from
distinct parallel lines; one must also test whether $P_2-P_1$ is parallel to the
common direction. Likewise, perpendicular directions do not make the geometric
lines perpendicular unless the lines also intersect.
\end{remarkbox}

\begin{derivationbox}[Testing intersection]
Solve
\[
P_1+t\mathbf v_1=P_2+s\mathbf v_2.
\]
This gives three scalar equations for two parameters. A common solution produces
the intersection point. If the directions are nonparallel but the system is
inconsistent, the lines are skew.
\end{derivationbox}

\begin{definitionbox}
Two lines in $\mathbb R^3$ are \emph{skew} when they are neither parallel nor
intersecting.
\end{definitionbox}

\begin{examplebox}
Consider
\[
L_1=(0,0,0)+t(1,0,0)
\]
and
\[
L_2=(0,1,1)+s(0,1,0).
\]
Their direction vectors have dot product zero. Yet every point on $L_1$ has
third coordinate $0$, while every point on $L_2$ has third coordinate $1$.
Thus the lines do not intersect: their directions are perpendicular, but the
lines are skew.
\end{examplebox}

\begin{center}
\begin{tikzpicture}[scale=0.9,>=stealth]
    \draw[blue!70!black,line width=1.1pt] (0.2,0.8) -- (5.4,0.8)
        node[right] {$L_1$};
    \draw[orange!85!black,line width=1.1pt] (2.0,2.0) -- (2.0,4.8)
        node[above] {$L_2$};
    \draw[gray,dashed] (2.0,2.0) -- (2.0,0.8);
    \node[gray,anchor=west] at (2.15,1.4) {separation in space};
    \draw[blue!70!black,->] (1.0,0.8) -- (2.2,0.8)
        node[midway,below] {$\mathbf v_1$};
    \draw[orange!85!black,->] (2.0,2.4) -- (2.0,3.6)
        node[midway,right] {$\mathbf v_2$};
\end{tikzpicture}
\end{center}

\begin{interpretationbox}[A genuinely three-dimensional possibility]
Skew lines have no analogue in the plane. In two dimensions, distinct
nonparallel lines must intersect. In space, one line can pass above or beside
another without meeting it.
\end{interpretationbox}

\begin{summarybox}[Lines in space]
A line is $P_0+t\mathbf v$. Two points determine the direction $Q-P$. A zero
cross product detects parallel directions, a zero dot product detects
perpendicular directions, and a parameter system distinguishes intersection
from skewness.
\end{summarybox}
